% ============================================================
% Mathematical Modeling Contest Paper Template (English)
%
% Compile: xelatex + bibtex
%   cd writer && xelatex main && bibtex main && xelatex main && xelatex main
%
% Structure: organized by sub-problem.
% 03-q1.tex is a skeleton for Problem 1.
% Writer creates 04-q2.tex, 05-q3.tex, etc. as needed.
% ============================================================
\documentclass[12pt,a4paper]{article}

% ---- Math ----
\usepackage{amsmath,amssymb,amsfonts}

% ---- Figures & Tables ----
\usepackage{graphicx}
\usepackage{float}
\usepackage{subfigure}
\usepackage{booktabs}
\usepackage{multirow}
\usepackage{longtable}

% ---- Bibliography ----
\usepackage[numbers,sort&compress]{natbib}
\bibliographystyle{plainnat}

% ---- Page Layout ----
\usepackage[left=2.5cm,right=2.5cm,top=2.5cm,bottom=2.5cm,headheight=14pt]{geometry}

% ---- Header & Footer ----
\usepackage{fancyhdr}
\pagestyle{fancy}
\fancyhf{}
\fancyhead[C]{\small Mathematical Modeling Contest Paper}
\fancyfoot[C]{\thepage}
\renewcommand{\headrulewidth}{0.4pt}

% ---- Code Listings ----
\usepackage{listings}
\usepackage{xcolor}
\lstset{
  basicstyle=\small\ttfamily,
  keywordstyle=\color{blue},
  commentstyle=\color{gray},
  stringstyle=\color{red!70!black},
  numbers=left,
  numberstyle=\tiny\color{gray},
  numbersep=8pt,
  frame=single,
  breaklines=true,
  breakatwhitespace=true,
  tabsize=4,
  showstringspaces=false,
  captionpos=b,
}

% ---- Hyperlinks ----
\usepackage[colorlinks=true,linkcolor=blue!70!black,citecolor=green!50!black,urlcolor=blue!80!black]{hyperref}

% ---- Other ----
\usepackage{enumitem}
\usepackage{caption}

% ---- Title ----
% Writer must replace this with a specific title based on the problem
\title{\textbf{(Set a specific title based on the problem)}}
\author{Team \#XXXXX}
\date{}

% ============================================================
\begin{document}

\maketitle
\thispagestyle{fancy}

% ============================================================
% Summary Sheet
% ============================================================
\begin{abstract}

(Write a one-page summary of your approach and key findings here.)

\vspace{1em}
\noindent\textbf{Keywords:} Keyword1, Keyword2, Keyword3
\end{abstract}


\newpage
\tableofcontents
\newpage

% === Body sections ===
% ============================================================
% Introduction
% ============================================================
\section{Introduction}

(Begin with the real-world context of the problem~\cite{example2024},
then describe the specific challenge posed by the contest.
Follow with a brief literature review of relevant approaches, citing 3-4 references.
Integrate the problem reformulation: translate the contest problem into precise
mathematical language -- define decision variables, objective functions, and constraints.
Conclude with an overview of the approach for each sub-problem.)

% ============================================================
% 2. Assumptions and Notation
% ============================================================
\section{Assumptions}

\begin{enumerate}
  \item Assumption 1: \ldots
  \item Assumption 2: \ldots
\end{enumerate}

\section{Notation}

\begin{table}[H]
\centering
\caption{Key Notation}
\begin{tabular}{cl}
\toprule
\textbf{Symbol} & \textbf{Description} \\
\midrule
$x$ & Example variable \\
$y$ & Example variable \\
\bottomrule
\end{tabular}
\end{table}


% === Per-problem chapters ===
% Each sub-problem gets its own file with: Model -> Solution -> Results
% ============================================================
% Problem 1: Modeling and Solution
% Writer creates independent files for each sub-problem (03-q1.tex, 04-q2.tex, ...)
% Each file contains the full structure: formulation -> solution -> results
% ============================================================
\section{Problem 1: Modeling and Solution}

\subsection{Problem Analysis}

(Analyze the core challenge and approach for Problem 1.
Use "we" as the subject: "For Problem 1, we observe that..."
Focus on: what is the key difficulty, and what is our high-level strategy?
Keep to 1-2 paragraphs.)

\subsection{Model Development}

(Start with the motivation: why this particular approach?
What phenomena or data characteristics led to this choice?
Then present the mathematical formulation with complete notation.
Connect equations with explanatory text -- avoid stacking formulas.)

\subsection{Model Solution}

(Describe the solution algorithm and computational process.
Place figures at the point of first reference using \verb|\begin{figure}[H]|.
Each figure/table should be followed by at least 3 sentences of analysis.)

% Figure example (Writer replaces with actual figures):
% \begin{figure}[H]
%   \centering
%   \includegraphics[width=0.8\textwidth]{figure/fig_q1_example.png}
%   \caption{Problem 1 result visualization}
%   \label{fig:q1_example}
% \end{figure}
%
% As shown in Figure~\ref{fig:q1_example}, we observe that...
% This indicates...
% Notably, ...

\subsection{Results and Discussion}

(Support every conclusion with specific numbers.
e.g., "The model achieves an $R^2$ of 0.943 and RMSE of 8.7, representing
a 12.8\% improvement over the baseline."
Include error analysis and sensitivity analysis where applicable.
Use tables for comparative results.)

% Results table example (Writer replaces with actual data):
% \begin{table}[H]
% \centering
% \caption{Problem 1: Model comparison}
% \begin{tabular}{lccc}
% \toprule
% Model & $R^2$ & RMSE & AIC \\
% \midrule
% Baseline & 0.813 & 15.2 & 2762 \\
% Proposed & 0.943 & 8.7 & 2408 \\
% \bottomrule
% \end{tabular}
% \end{table}

% Writer creates additional files as needed:
% \input{04-q2}
% \input{05-q3}

% === Cross-cutting analysis ===
% ============================================================
% Sensitivity Analysis
% ============================================================
\section{Sensitivity Analysis}

(Select 2-4 key parameters in the model and systematically vary each
while holding others constant. Present results as line plots showing how
the output metric changes with each parameter.

Address: Which parameters most influence the result? Is the model robust
to reasonable perturbations? Are there parameter ranges where the model
breaks down? Discuss real-world implications.)

% ============================================================
% Strengths and Weaknesses
% ============================================================
\section{Strengths and Weaknesses}

\subsection{Strengths}

(Discuss in paragraph form, not bullet lists. Support each strength with
evidence from the results: accuracy metrics, efficiency, interpretability,
generalizability. Compare against baseline or alternative approaches.)

\subsection{Weaknesses and Future Work}

(Honestly discuss limitations in paragraph form. Does the model depend on
assumptions that may not hold? Is it sensitive to data quality? Each
weakness should be paired with a potential mitigation or future direction.)

% ============================================================
% Conclusions
% ============================================================
\section{Conclusions}

(Concisely restate the key findings for each sub-problem with quantitative
results. Highlight the most novel aspect of the approach. End with 2-3
sentences on broader implications or future extensions.
Do not introduce new results or methods here.)


% === References ===
\bibliography{references}

% === Appendix ===
\appendix
% ============================================================
% Appendix
% ============================================================
\section{Appendix}

\subsection{Core Code}

% Writer must use \lstinputlisting to include Coder's actual code files.
% Path is relative to writer/ directory, using ../coder/ prefix.
% Each code segment should have a brief text introduction.
% Do NOT use \begin{lstlisting} to manually type code.

% The following implements the core prediction model for Problem 1:
% \lstinputlisting[language=Python, caption={Problem 1: Prediction Model},
%   firstline=10, lastline=80]{../coder/solve_q1.py}

\subsection{Supplementary Figures}

% Additional figures that support the analysis but would clutter the main text.

\subsection{Detailed Derivations}

% Lengthy mathematical derivations summarized in the body.


\end{document}
